\section*{Pregunta 2}
\noindent El primer prototipo debe ser diseñada como un sólido que se forma en el primer octante con vértices en \(P(0,0,0)\), \(A(13,0,0)\), \(B(0,11,0)\) y \(C(0,0,16)\)

\begin{itemize}
    \item[a.] Represente una función de varias variables que modele la superficie inclinada del sólido.

    \item[c.] Mediante el uso de las la integral doble calcule el volumen del sólido.
\end{itemize}

\subsection*{a.- Represente una función de varias variables que modele la superficie inclinada del sólido}

\begin{multicols}{2}
\begin{center}
\begin{tabular}{|c|c|c|c|}
    \hline
	& x & y & z \\
	\hline
	Eje x & 13 & 0 & 0 \\
	\hline
	Eje y & 0 & 11 & 0 \\
	\hline
	Eje z & 0 & 0 & 16 \\
    \hline
\end{tabular}
\end{center}

\columnbreak

\begin{align*}
	1 & = \frac{x}{13} + \frac{y}{11} + \frac{z}{16} \\
	z & = 16 \left(1 - \frac{x}{13} - \frac{y}{11}\right)
\end{align*}
\end{multicols}

\section*{Mediante el uso de las la integral doble calcule el volumen del sólido}

\begin{gather*}
	0 = 1 - \frac{x}{13} - \frac{y}{11} \\
	0 = \frac{13 - x}{13} - \frac{y}{11} \\
	\frac{y}{11} = \frac{13 - x}{13} \\
	y = 11 \left(\frac{13 - x}{13}\right) \\
	y = \frac{11(13 - x)}{13} \\
	y = \frac{143 - 11x}{13} \\
	\rule{0.5\textwidth}{0.4pt} \\
	\int_{0}^{13} \int_{0}^{\frac{143-11x}{13}} 16 \left(1 - \frac{x}{13} - \frac{y}{11}\right) dy dx \\
	\frac{16}{143} \int_{0}^{13} \int_{0}^{\frac{143-11x}{13}} (143 - 11x - 13y) dy dx \\
	\frac{16}{143} \int_{0}^{13} \left[ 143y - 11xy - \frac{13y^2}{2} \right]_{0}^{\frac{143-11x}{13}} dx \\
	\frac{16}{143} \int_{0}^{13} -7 \cdot 143^2 + 26 \cdot 1573x - 7 \cdot 121x^2 dx \\
	\frac{16}{143} \left[ -7 \cdot 143^2x + \frac{26 \cdot 1573x^2}{2} - \frac{7 \cdot 121 x^3}{3} \right]_{0}^{13} \\
	\frac{327124}{3} = 109061,\bar{3}
\end{gather*}